\section{Curvature}

\vspace{0.4cm}\emph{\textbf{(Lecture 9)}}\\


We will study the curvature of a Riemannian manifold from an intrinsic viewpoint: relative acceleration of nearby geodesics, angles of triangles, etc.

\begin{colordef}{Riemann curvature tensor}{riemann-curvature-tensor}
The \emph{Riemann curvature tensor} $R\in T^1_3(M)$ of a connection $\nabla$ is defined as
\[
R(X,Y)\,Z = \nabla_X \nabla_Y Z - \nabla_Y \nabla_X Z - \nabla_{[X,Y]} Z.
\]
\end{colordef}

One can view $R(X,Y)$ as an operator from $\Gamma(TM)$ to itself via $R(X,Y) = [\nabla_X, \nabla_Y] - \nabla_{[X,Y]}$, so $R(X,Y)\,Z\in \Gamma(TM)$.

\begin{colortheo}{$R$ is a tensor field}{curvature_tensor_is_tensor}
$R$ is $C^\infty$-linear in all its arguments.
\end{colortheo}

\begin{proof}
It suffices to prove it is $C^\infty$-linear in all its arguments.

Recall: $[fX, Y] = f[X,Y] - Y(f)\,X$.

So
\begin{align*}
R(fX, Y)\,Z &= \nabla_{fX} \nabla_Y Z - \nabla_Y \nabla_{fX} Z - \nabla_{[fX,Y]} Z \\
&= f\,\nabla_X \nabla_Y Z - \nabla_Y(f\,\nabla_X Z) - \nabla_{f[X,Y] - Y(f)X} Z \\
&= f\,\nabla_X \nabla_Y Z - Y(f)\,\nabla_X Z - f\,\nabla_Y \nabla_X Z - f\,\nabla_{[X,Y]} Z + Y(f)\,\nabla_X Z \\
&= f\,R(X,Y)\,Z.
\end{align*}
Similarly for the other arguments.
\end{proof}

We will be only interested in the Riemann curvature tensor of the Levi-Civita connection of $(M,g)$.

To get some geometrical intuition of $R$ consider the following situation.


Let $X, Y$ be coordinate vector fields, so $[X,Y] = 0$, by \eqref{eq:coordinate_fields_commute}. Starting from $p \in M$ and successively flowing along $X$, $Y$, $-X$, $-Y$ for small times $s_1, s_2, s_1, s_2$ respectively, we obtain a closed loop $\gamma(s_1, s_2)$ at $p$ (the loop closes exactly because the flows of $X$ and $Y$ commute). Let $P \colon T_pM \to T_pM$ denote parallel transport once around this loop. Then we have the following

\begin{colorlem}{Characterisation of the Riemann curvature tensor}{}
\[
\frac{P Z - Z}{s_1 \cdot s_2} \xrightarrow{s_1, s_2 \to 0} R(X,Y) \cdot Z.
\]
\end{colorlem}
\begin{proof}
See exercise 9.4b.
\end{proof}


% \begin{center}
% \begin{tikzpicture}[scale=4, >=stealth, line join=round]
%   % the loop
%   \draw[->, thick] (0,0) -- (1,0);
%   \draw[->, thick] (1,0) -- (1,1);
%   \draw[->, thick] (1,1) -- (0,1);
%   \draw[->, thick] (0,1) -- (0,0);
 
%   % corner labels
%   \node[below left]  at (0,0) {$p$};
 
%   % edge labels
%   \node[below] at (0.5, 0)   {$s_1\, X$};
%   \node[right] at (1, 0.5)   {$s_2\, Y$};
%   \node[above] at (0.5, 1)   {$-s_1\, X$};
%   \node[left]  at (0, 0.5)   {$-s_2\, Y$};
 
%   % the vector Z at p, and its image after transport
%   \draw[->, blue, thick] (0,0) -- (0.35, 0.5);
%   \node[blue, right] at (0.32, 0.45) {$Z$};
 
%   \draw[->, red!70!black, thick, dashed] (0,0) -- (0.5, 0.35);
%   \node[red!70!black, below right] at (0.45, 0.30) {$P Z$};
% \end{tikzpicture}
% \end{center}

\begin{figure}[H]
\begin{center}
\begin{tikzpicture}[scale=5, >=stealth, line join=round]
  % the loop
  \draw[->, thick] (0,0) -- (1,0) node[midway, below] {$s_1 X$};
  \draw[->, thick] (1,0) -- (1,1) node[midway, right] {$s_2 Y$};
  \draw[->, thick] (1,1) -- (0,1) node[midway, above] {$-s_1 X$};
  \draw[->, thick] (0,1) -- (0,0) node[midway, left]  {$-s_2 Y$};

  \node[below left] at (0,0) {$p$};

  % Z before transport
  \draw[->, Canard, very thick] (0,0) -- (0.30, 0.55) node[pos=1, above] {$Z$};

  % PZ after transport (slightly rotated to show discrepancy)
  \draw[->, Groseille, very thick, dashed] (0,0) -- (0.50, 0.40) node [pos=1, below right] {$PZ$};

  % the gap
  \draw[->, Ardoise] (0.30, 0.55) -- (0.50, 0.40);
  \node[Ardoise] at (0.7, 0.5) {$\sim s_1 s_2 R(X,Y)Z$};
\end{tikzpicture}
\end{center}
\end{figure}


So $R(X,Y)Z$ measures, to leading order, the failure of parallel transport around an infinitesimal coordinate loop of area $s_1 s_2$ to return $Z$ to itself.\\
 
\begin{colremark}
This gives a heuristic explanation for one of the symmetries of $R$ that we will prove rigorously in Theorem~\ref{thm:symmetries-riemann-curvature-tensor}. Since $P$ is a linear isometry, expanding $\langle PZ, PW\rangle = \langle Z, W\rangle$ to leading order in $s_1 s_2$ forces
\[
\langle R(X,Y)Z, W\rangle + \langle Z, R(X,Y)W\rangle = 0,
\]
i.e.\ $R(X,Y)$ is skew-symmetric in its last two arguments.
\end{colremark}


Recall the musical isomorphism $\flat \colon TM \to T^*M$, $X \mapsto X^\flat = g(X, \cdot)$, with components $X_a = g_{ab} X^b$. We use it to lower the vector-valued output of $R$ to a $(0,4)$-tensor. (Since $\nabla g = 0$ for the Levi-Civita connection, lowering and raising indices commutes with $\nabla$, which will be convenient later.)



% \begin{colordef}{The (0,4) Riemann curvature tensor}{04-riemann-curvature-tensor}
% We define the $(0,4)$-Riemann curvature tensor $\widehat{R} \in T^0_4(M)$ by
% \begin{align*}
%     \widehat{R}(X,Y,Z,W) &:= g(R(X,Y)\cdot W,\, Z) \\
%     &= -\,g(R(X,Y)\cdot Z,\, W)\\
%     &= (R(X,Y)W)^\flat(Z).
% \end{align*}
% \end{colordef}

% Be aware, that this tensor is often the subject of confusion: different authors have different sign conventions, and further, the $\widehat{\cdot}$ symbol is often not written explicitly. We will also adopt this convention. The type of arguments the tensor is evaluated on will make clear which version of the tensor we are using.

% The components of $R$ as a $(1,3)$-tensor are defined by
% \[
% R(\partial_\gamma, \partial_\delta)\,\partial_\beta = R^\alpha{}_{\beta\gamma\delta}\,\partial_\alpha,
% \]
% so that $R^\alpha{}_{\beta\gamma\delta} = \big(R(\partial_\gamma, \partial_\delta)\,\partial_\beta\big)^\alpha$. The components of the $(0,4)$-tensor $\widehat R$ are defined as usual by evaluation on the coordinate basis,
% \[
% \widehat R_{\alpha\beta\gamma\delta} = \widehat R(\partial_\alpha, \partial_\beta, \partial_\gamma, \partial_\delta) = g(\partial_\alpha,\, R(\partial_\gamma, \partial_\delta)\,\partial_\beta) = g_{\alpha\mu}\, R^\mu{}_{\beta\gamma\delta}.
% \]
% From now on we drop the hat and write
% \begin{equation}
%     \label{eq:riemann-tensor-notation}
%     R_{\alpha\beta\gamma\delta} := g_{\alpha\mu}\, R^\mu{}_{\beta\gamma\delta}.
% \end{equation}



\begin{colordef}{The (0,4) Riemann curvature tensor}{04-riemann-curvature-tensor}
We define the $(0,4)$-Riemann curvature tensor $\widehat{R} \in T^0_4(M)$ by
\[
\widehat R(X, Y, Z, W) := g\big(X,\, R(Z, W)\, Y\big).
\]
Equivalently, using the musical isomorphism $\flat$:
\[
\widehat R(X, Y, Z, W) = X^\flat\big(R(Z, W)\, Y\big).
\]
\end{colordef}

Be aware, that this tensor is often the subject of confusion: different authors have different sign conventions, and further, the $\widehat{\cdot}$ symbol is often not written explicitly. We will also adopt this convention. The type of arguments the tensor is evaluated on will make clear which version of the tensor we are using.

The components of $R$ as a $(1,3)$-tensor are defined by
\[
R(\partial_\gamma, \partial_\delta)\,\partial_\beta = R^\alpha{}_{\beta\gamma\delta}\,\partial_\alpha,
\]
so that $R^\alpha{}_{\beta\gamma\delta} = \big(R(\partial_\gamma, \partial_\delta)\,\partial_\beta\big)^\alpha$. The components of the $(0,4)$-tensor $\widehat R$ are defined as usual by evaluation on the coordinate basis,
\[
\widehat R_{\alpha\beta\gamma\delta} = \widehat R(\partial_\alpha, \partial_\beta, \partial_\gamma, \partial_\delta) = g\big(\partial_\alpha,\, R(\partial_\gamma, \partial_\delta)\,\partial_\beta\big) = g_{\alpha\mu}\, R^\mu{}_{\beta\gamma\delta}.
\]
From now on we drop the hat and write
\begin{equation}
    \label{eq:riemann-tensor-notation}
    R_{\alpha\beta\gamma\delta} := g_{\alpha\mu}\, R^\mu{}_{\beta\gamma\delta}.
\end{equation}











% \begin{colorlem}{Components of the Riemann tensor}{}
% \[
% R^a{}_{b\gamma\delta} = \partial_\gamma \Gamma^a_{\delta b} - \partial_\delta \Gamma^a_{\gamma b} + \Gamma^a_{\gamma \lambda}\,\Gamma^\lambda_{\delta b} - \Gamma^a_{\delta \lambda}\,\Gamma^\lambda_{\gamma b},
% \]
% and
% \[
% R_{\alpha\beta\gamma\delta} = \frac{1}{2}\left(\partial^2_{\beta\gamma}\,g_{\alpha\delta} - \partial^2_{\alpha\gamma}\,g_{\beta\delta} + \partial^2_{\alpha\delta}\,g_{\beta\gamma} - \partial^2_{\beta\delta}\,g_{\alpha\gamma}\right) + \partial g \cdot \partial g.
% \]
% \end{colorlem}


% \begin{proof}
%     old one
% \begin{align*}
% R^a{}_{b\gamma\delta} &= \big(R(\partial_\gamma, \partial_\delta)\,\partial_b\big)^a \\
% &= \big(\nabla_{\partial_\gamma} \nabla_{\partial_\delta} \partial_b - \nabla_{\partial_\delta} \nabla_{\partial_\gamma} \partial_b\big)^a \\
% &= \partial_\gamma (\nabla_{\partial_\delta} \partial_b)^a + \Gamma^a_{\gamma\lambda}\,(\nabla_{\partial_\delta} \partial_b)^\lambda - (\gamma \leftrightarrow \delta) \\
% &= \partial_\gamma\!\left(\partial_\delta (\partial_b)^a + \Gamma^a_{\delta k}\,(\partial_b)^k\right) + \Gamma^a_{\gamma\lambda}\!\left(\partial_\delta(\partial_b)^\lambda + \Gamma^\lambda_{\delta k}\,(\partial_b)^k\right) - (\gamma \leftrightarrow \delta) \\
% &= \partial_\gamma \Gamma^a_{\delta b} + \Gamma^a_{\gamma\lambda}\,\Gamma^\lambda_{\delta b} - (\gamma \leftrightarrow \delta). \qedhere
% \end{align*}
% \end{proof}





\begin{colorlem}{Components of the Riemann tensor}{}
In local coordinates, the components of the Riemann tensor are
\begin{equation}
    \label{eq:riemann-curvature-tensor-components1}
    R^a{}_{b\gamma\delta} = \partial_\gamma \Gamma^a_{\delta b} - \partial_\delta \Gamma^a_{\gamma b} + \Gamma^a_{\gamma\lambda}\,\Gamma^\lambda_{\delta b} - \Gamma^a_{\delta\lambda}\,\Gamma^\lambda_{\gamma b}.
\end{equation}
The components of the lowered tensor $R_{\alpha\beta\gamma\delta} = g_{\alpha\mu}\,R^\mu{}_{\beta\gamma\delta}$ can be written as
\begin{equation}
    \label{eq:riemann-curvature-tensor-components2}
    R_{\alpha\beta\gamma\delta} = \tfrac{1}{2}\left(\partial^2_{\beta\gamma}\,g_{\alpha\delta} - \partial^2_{\alpha\gamma}\,g_{\beta\delta} + \partial^2_{\alpha\delta}\,g_{\beta\gamma} - \partial^2_{\beta\delta}\,g_{\alpha\gamma}\right) + Q(g, \partial g),
\end{equation}
where $Q(g, \partial g)$ is a function of $g$, its inverse $g^{-1}$, and first derivatives $\partial g$, but no higher derivatives.
\end{colorlem}


We will derive the first formula here. Equation \eqref{eq:riemann-curvature-tensor-components2} is a consequence of the first formula and the Christoffel formula, Definition \ref{def:christoffel}.

\begin{proof}

By the definition of the Riemann tensor (Definition~\ref{def:riemann-curvature-tensor}) and the fact that coordinate vector fields commute, \eqref{eq:coordinate_fields_commute},
\[
R(\partial_\gamma, \partial_\delta)\,\partial_b = \nabla_{\partial_\gamma} \nabla_{\partial_\delta} \partial_b - \nabla_{\partial_\delta} \nabla_{\partial_\gamma} \partial_b.
\]

By the definition of the Christoffel symbols (Definition~\ref{def:christoffel}):
\[
\nabla_{\partial_\delta} \partial_b = \Gamma^\lambda_{\delta b}\,\partial_\lambda.
\]

Applying $\nabla_{\partial_\gamma}$ to both sides and using the Leibniz rule for the connection (Definition~\ref{def:connection}, axiom 3):
\begin{align*}
\nabla_{\partial_\gamma}\!\left(\Gamma^\lambda_{\delta b}\,\partial_\lambda\right) 
&= \partial_\gamma\!\left(\Gamma^\lambda_{\delta b}\right)\,\partial_\lambda + \Gamma^\lambda_{\delta b}\,\nabla_{\partial_\gamma}\partial_\lambda \\
&= \partial_\gamma\!\left(\Gamma^\lambda_{\delta b}\right)\,\partial_\lambda + \Gamma^\lambda_{\delta b}\,\Gamma^a_{\gamma\lambda}\,\partial_a.
\end{align*}

Renaming the dummy index $\lambda \to a$ in the first term:
\[
\nabla_{\partial_\gamma}\!\left(\Gamma^\lambda_{\delta b}\,\partial_\lambda\right) = \partial_\gamma\!\left(\Gamma^a_{\delta b}\right)\,\partial_a + \Gamma^a_{\gamma\lambda}\,\Gamma^\lambda_{\delta b}\,\partial_a.
\]

The analogous expression with $\gamma$ and $\delta$ swapped is:
\[
\nabla_{\partial_\delta}\!\left(\Gamma^\lambda_{\gamma b}\,\partial_\lambda\right) = \partial_\delta\!\left(\Gamma^a_{\gamma b}\right)\,\partial_a + \Gamma^a_{\delta\lambda}\,\Gamma^\lambda_{\gamma b}\,\partial_a.
\]

Subtracting:
\[
R(\partial_\gamma, \partial_\delta)\,\partial_b = \left(\partial_\gamma \Gamma^a_{\delta b} - \partial_\delta \Gamma^a_{\gamma b} + \Gamma^a_{\gamma\lambda}\,\Gamma^\lambda_{\delta b} - \Gamma^a_{\delta\lambda}\,\Gamma^\lambda_{\gamma b}\right)\partial_a.
\]

Reading off the $\partial_a$-component yields \eqref{eq:riemann-curvature-tensor-components1}.

% \textbf{Step 2: derivation of $R_{\alpha\beta\gamma\delta}$.}

% Lowering the index using $g$:
% \[
% R_{\alpha\beta\gamma\delta} = g_{\alpha\mu}\,R^\mu{}_{\beta\gamma\delta} = g_{\alpha\mu}\!\left(\partial_\gamma \Gamma^\mu_{\delta\beta} - \partial_\delta \Gamma^\mu_{\gamma\beta} + \Gamma^\mu_{\gamma\lambda}\,\Gamma^\lambda_{\delta\beta} - \Gamma^\mu_{\delta\lambda}\,\Gamma^\lambda_{\gamma\beta}\right).
% \]

% The last two terms are products of Christoffel symbols, hence (by the formula for Christoffel symbols in terms of $g$, equation~\eqref{eq:christoffel-formula}) products of $g^{-1}$ and $\partial g$ — these contribute to $Q(g, \partial g)$.

% For the first two terms, recall the Christoffel formula:
% \[
% \Gamma^\mu_{\delta\beta} = \tfrac{1}{2}\,g^{\mu\nu}\!\left(\partial_\delta g_{\beta\nu} + \partial_\beta g_{\delta\nu} - \partial_\nu g_{\delta\beta}\right).
% \]

% Therefore:
% \[
% g_{\alpha\mu}\,\partial_\gamma \Gamma^\mu_{\delta\beta} = g_{\alpha\mu}\,\partial_\gamma\!\left[\tfrac{1}{2}\,g^{\mu\nu}\!\left(\partial_\delta g_{\beta\nu} + \partial_\beta g_{\delta\nu} - \partial_\nu g_{\delta\beta}\right)\right].
% \]

% Applying the product rule, the only second-derivative-of-$g$ terms come from $\partial_\gamma$ hitting the inner parenthesis (not $g^{\mu\nu}$ or $g_{\alpha\mu}$, since those only contribute first derivatives). The contraction $g_{\alpha\mu}\,g^{\mu\nu} = \delta_\alpha^\nu$ then collapses the $\nu$ index to $\alpha$:
% \[
% g_{\alpha\mu}\,\partial_\gamma \Gamma^\mu_{\delta\beta} = \tfrac{1}{2}\!\left(\partial^2_{\gamma\delta} g_{\beta\alpha} + \partial^2_{\gamma\beta} g_{\delta\alpha} - \partial^2_{\gamma\alpha} g_{\delta\beta}\right) + (\text{terms with } \partial g \cdot \partial g).
% \]

% Similarly:
% \[
% g_{\alpha\mu}\,\partial_\delta \Gamma^\mu_{\gamma\beta} = \tfrac{1}{2}\!\left(\partial^2_{\delta\gamma} g_{\beta\alpha} + \partial^2_{\delta\beta} g_{\gamma\alpha} - \partial^2_{\delta\alpha} g_{\gamma\beta}\right) + (\text{terms with } \partial g \cdot \partial g).
% \]

% Subtracting, the $\partial^2_{\gamma\delta}\,g_{\beta\alpha}$ and $\partial^2_{\delta\gamma}\,g_{\beta\alpha}$ terms cancel (partial derivatives commute), leaving:
% \[
% g_{\alpha\mu}\!\left(\partial_\gamma \Gamma^\mu_{\delta\beta} - \partial_\delta \Gamma^\mu_{\gamma\beta}\right) = \tfrac{1}{2}\!\left(\partial^2_{\gamma\beta} g_{\delta\alpha} - \partial^2_{\gamma\alpha} g_{\delta\beta} - \partial^2_{\delta\beta} g_{\gamma\alpha} + \partial^2_{\delta\alpha} g_{\gamma\beta}\right) + Q(g, \partial g).
% \]

% Relabelling to match the stated formula (using symmetry $\partial^2_{\gamma\beta} = \partial^2_{\beta\gamma}$):
% \[
% R_{\alpha\beta\gamma\delta} = \tfrac{1}{2}\!\left(\partial^2_{\beta\gamma}\,g_{\alpha\delta} - \partial^2_{\alpha\gamma}\,g_{\beta\delta} + \partial^2_{\alpha\delta}\,g_{\beta\gamma} - \partial^2_{\beta\delta}\,g_{\alpha\gamma}\right) + Q(g, \partial g). \qedhere
% \]
\end{proof}


\begin{colremark}{}{}
Note: $R = 0$ on $(\mathbb{R}^n, g_E)$. Conversely, one can show that $(M,g)$ is locally isometric to $(\mathbb{R}^n, g_E)$ if and only if $R = 0$.
\end{colremark}




% \begin{colremark}{}{}
% In normal coordinates around $p$: $g|_p = \delta_{ij}$, $\partial g|_p = 0$, see Lemma \ref{lem:metric_normal_coordinates}. So there, we have
% \begin{equation}
%     \label{eq:riemann-curvature-normal-coordinates}
%     R_{\alpha\beta\gamma\delta}\big|_p = \frac{1}{2}\left(\partial^2_{\beta\gamma}\,g_{\alpha\delta} - \partial^2_{\alpha\gamma}\,g_{\beta\delta} + \partial^2_{\alpha\delta}\,g_{\beta\gamma} - \partial^2_{\beta\delta}\,g_{\alpha\gamma}\right)\bigg|_p.
% \end{equation}

% In addition, by the Gauss lemma: in normal coordinates \eqref{eq:gauss_lemma_radial_vector_field} holds not just at $p$, i.e. $g_{ij}\,x^i = \delta_{ij}\,x^i$. Differentiating yields
% \[
% \left(\partial_\gamma\,g_{\beta\delta} + \partial^2_{\gamma\alpha}\,g_{\beta\delta} + \partial^2_{\beta\gamma}\,g_{\alpha\delta}\right)\big|_p = 0.
% \]

% So in normal coordinates at $p$:
% \[
% R_{\alpha\beta\gamma\delta}\big|_p = \left(\partial^2_{\beta\gamma}\,g_{\alpha\delta} - \partial^2_{\alpha\gamma}\,g_{\beta\delta}\right)\big|_p.
% \]
% \end{colremark}








\begin{colremark}{}{}
In normal coordinates around $p$, by Lemma~\ref{lem:metric_normal_coordinates}, $g_{ij}|_p = \delta_{ij}$ and $\partial g|_p = 0$. So the formula for $R_{\alpha\beta\gamma\delta}$ from \eqref{eq:riemann-curvature-tensor-components2} simplifies at $p$ to
\begin{equation}
    \label{eq:riemann-curvature-normal-coordinates}
    R_{\alpha\beta\gamma\delta}\big|_p = \tfrac{1}{2}\!\left(\partial^2_{\beta\gamma}\,g_{\alpha\delta} - \partial^2_{\alpha\gamma}\,g_{\beta\delta} + \partial^2_{\alpha\delta}\,g_{\beta\gamma} - \partial^2_{\beta\delta}\,g_{\alpha\gamma}\right)\bigg|_p,
\end{equation}
since the derivative terms $\partial g$ terms vanish at $p$. Moreover, by the Gauss lemma, the identity \eqref{eq:gauss_lemma_radial_vector_field} holds on the entire normal neighborhood (not just at $p$):
\[
g_{ij}(x)\,x^i = x_j.
\]
Differentiating three times with respect to $x^\delta$, $x^\beta$, $x^\gamma$ and evaluating at $p$ (where $x = 0$ and $\partial g|_p = 0$) yields the cyclic identity
\[
\partial^2_{\beta\gamma}\,g_{\alpha\delta}\big|_p + \partial^2_{\gamma\delta}\,g_{\alpha\beta}\big|_p + \partial^2_{\delta\beta}\,g_{\alpha\gamma}\big|_p = 0.
\]
Substituting this into \eqref{eq:riemann-curvature-normal-coordinates}, the formula simplifies to
\[
R_{\alpha\beta\gamma\delta}\big|_p = \left(\partial^2_{\beta\gamma}\,g_{\alpha\delta} - \partial^2_{\alpha\gamma}\,g_{\beta\delta}\right)\big|_p.
\]
\end{colremark}



\begin{colortheo}{Symmetries of the Riemann curvature tensor}{symmetries-riemann-curvature-tensor}
\begin{enumerate}
\item Antisymmetry in the first pair: $R(X,Y)\,Z = -R(Y,X)\,Z$
\item Antisymmetry in the last arguments: $\langle R(X,Y)\cdot W, Z \rangle = -\langle R(X,Y)\cdot Z, W \rangle$
\item Symmetry pairwise: ${R}(X,Y,Z,W) = {R}(Z,W,X,Y)$
\item $1^\text{st}$ Bianchi identity: $R(X,Y,Z,W) + R(Z,X,Y,W) + R(Y,Z,X,W) = 0$
\item $2^\text{nd}$ Bianchi identity: $(\nabla_X R)(Y,Z,V,W) + (\nabla_Z R)(X,Y,V,W) + (\nabla_Y R)(Z,X,V,W) = 0$
\end{enumerate}
\end{colortheo}

Recall \eqref{eq:riemann-tensor-notation} for the notation of $R_{\alpha\beta\gamma\delta}$: the last three bullets are the $(0,4)$-Riemann curvature tensor, and the first two are the $(1,3)$-Riemann curvature tensor.

\begin{proof}
It suffices to prove the identities in one coordinate system for coordinate vector fields. We choose normal coordinates at $p$. The statements become:

\begin{enumerate}
    \item[i-iii)] $R_{\alpha\beta\gamma\delta} = -R_{\beta\alpha\gamma\delta} = -R_{\alpha\beta\delta\gamma} = R_{\gamma\delta\alpha\beta}$
    
    \item[iv)] $R_{\alpha\beta\gamma\delta} + R_{\gamma\alpha\beta\delta} + R_{\beta\gamma\alpha\delta} = 0$
    
    \item[v)] $\nabla_\alpha R_{\beta\gamma\delta\varepsilon} + \nabla_\gamma R_{\alpha\beta\delta\varepsilon} + \nabla_\beta R_{\gamma\alpha\delta\varepsilon} = 0$
\end{enumerate}

Identities i-iii) follow immediately from the explicit expression \eqref{eq:riemann-curvature-normal-coordinates}, and iv) follows from i-iii). For v), since $\Gamma|_p = 0$, we have
\[
    \nabla_\alpha R_{\beta\gamma\delta\varepsilon}\big|_p = \partial_\alpha R_{\beta\gamma\delta\varepsilon}\big|_p,
\]
and the cyclic sum of $\partial_\alpha R_{\beta\gamma\delta\varepsilon}$ at $p$ vanishes by direct computation from \eqref{eq:riemann-curvature-normal-coordinates}.
\end{proof}


\begin{colremark}[$R$ as a map on bivectors]{}
\label{rem:R-map-on-bivectors111}
Since $R$ is antisymmetric in its first two arguments, we can view it as a map on bivectors:
\[
R(X, Y) = \widetilde R(X \wedge Y),
\]
where $X \wedge Y \in \Lambda^2 T_p M$. (Recall from Section \ref{sec:notation-tensors} that $\wedge$ can be applied to vectors as well as differential forms.) The antisymmetry $X \wedge Y = -Y \wedge X$ in $\Lambda^2 T_p M$ then automatically encodes the antisymmetry $R(X, Y) = -R(Y, X)$.
This is a special case of the following general algebraic fact. Let $V,W$ be vector spaces and $B: V \times V \to W$ a bilinear, alternating map (i.e.\ $B(X, Y) = -B(Y, X)$ and $B(X, X) = 0$ for any $X$). Then $B$ \emph{factors through} $\wedge: V\times V\to \Lambda^2 V$: there is a unique linear map $\widetilde B: \Lambda^2 V \to W$ such that
\[
B(X, Y) = \widetilde B(X \wedge Y) \qquad \text{for all } X, Y \in V.
\]
One can check this by defining $\widetilde B(X \wedge Y) := B(X, Y)$ on simple bivectors and extending by linearity, then verifying well-definedness (the relations defining $\Lambda^2 V$ are respected) and uniqueness (simple bivectors $X \wedge Y$ span $\Lambda^2 V$).
\[
\begin{tikzcd}[row sep=large, column sep=large]
V \times V \arrow[r, "\wedge"] \arrow[dr, "B"'] & \Lambda^2 V \arrow[d, dashed, "\widetilde B"] \\
& W
\end{tikzcd}
\]
\end{colremark}

\begin{colorlem}{Geodesic deviation equation}{geodesic-deviation}
Let $(M, g)$ be a smooth Riemannian manifold and let $\phi : (-\epsilon, \epsilon) \times [0, 1] \to M$ be a smooth map such that, for each $s \in (-\epsilon, \epsilon)$, the curve $\gamma_s = \phi(s, \cdot)$ is a geodesic. Define the vector fields along $\phi$
\[
T = \mathrm{d}\phi\!\left(\frac{\partial}{\partial t}\right), \qquad X = \mathrm{d}\phi\!\left(\frac{\partial}{\partial s}\right).
\]
Then
\[
\nabla_T \nabla_T X = -R(X, T)\, T.
\]
\end{colorlem}

\begin{proof}
See exercise 9.3.
\end{proof}

Here, intuitively, $X$ measures the infinitesimal separation between nearby geodesics; thus, the Riemann curvature tensor measures the \emph{relative acceleration} of nearby geodesics.







\begin{colordef}{Flat manifold}{flat-manifold}
A Riemannian manifold $(M, g)$ is called \emph{(locally) flat} if every point $p \in M$ has a neighborhood that is isometric to an open subset of Euclidean space $(\mathbb{R}^n, g_E)$. Equivalently, around each $p$ one can find coordinates in which $g_{ij} = \delta_{ij}$ on a neighborhood of $p$.
\end{colordef}

\begin{colremark}{}{}
This is a \emph{local} condition. Globally, a flat manifold need not look like $\mathbb{R}^n$. The flat torus $T^n = \mathbb{R}^n / \mathbb{Z}^n$ and the flat cylinder $S^1 \times \mathbb{R}$ are both locally flat but globally distinct from $\mathbb{R}^n$.
\end{colremark}

\begin{colorlem}{Flat $\iff$ $R = 0$}{flat-iff-R-zero}
A Riemannian manifold $(M, g)$ is locally flat if and only if $R \equiv 0$.
\end{colorlem}

\begin{proof}[Proof sketch of the easy direction]
If $(M, g)$ is locally flat, then in coordinates where $g_{ij} = \delta_{ij}$ all Christoffel symbols vanish, so $R = 0$ by the formula \eqref{eq:riemann-curvature-tensor-components1}. Since $R$ is a tensor, (Theorem \ref{thm:curvature_tensor_is_tensor}) this holds in every coordinate system.

The other direction is more complicated.
\end{proof}



\begin{colremark}[Degrees of freedom of the curvature tensor]
In 2D, there is only one independent component. This single number is the \emph{Gaussian curvature}. 

Using a combinatorial argument, one can show that in higher dimensions, we have more degrees of freedom. One can show that in 3D, we have 6 independent components. In 4D, we have 20 independent components.

\end{colremark}










\subsection{Sectional curvature}


\begin{colordef}{Sectional curvature}{sectional-curvature}
Let $(M, g)$ be a Riemannian manifold and $\Pi \subseteq T_p M$ a 2-dimensional subspace spanned by linearly independent vectors $X, Y \in T_p M$. The \emph{sectional curvature} associated to $\Pi$ is
\[
K(X, Y) := \frac{R(X, Y, X, Y)}{\|X\|^2\,\|Y\|^2 - \langle X, Y\rangle^2}.
\]
\end{colordef}

\begin{colremark}{}{}
The numerator $R(X, Y, X, Y)$ is the $(0,4)$-tensor evaluated on $(X, Y, X, Y)$; equivalently, in $(1,3)$-form,
\[
R(X, Y, X, Y) = \langle R(X, Y)\, Y,\, X\rangle.
\]
The denominator $\|X\|^2 \|Y\|^2 - \langle X, Y\rangle^2 = \|X \wedge Y\|^2$ is the squared area of the parallelogram spanned by $X$ and $Y$, so $K(\Pi)$ is a \emph{ratio of areas} and depends only on the plane $\Pi$, not on the choice of spanning vectors. We will verify this independence in the next lemma.
\end{colremark}



\begin{colorlem}{Sectional curvature}{}
$K(X,Y)$ is independent of the choice of $X, Y$ spanning $\Pi$.
\end{colorlem}

\begin{colremark}{}{}
As $K$ depends only on the plane $\Pi$, we write also $K(\Pi)$ instead of $K(X,Y)$.
\end{colremark}

\begin{proof}
Let $e_1, e_2$ be an orthonormal basis of $\Pi$, so that $X = X^1 e_1 + X^2 e_2$ and $Y = Y^1 e_1 + Y^2 e_2$.

Then due to the symmetries of $R$ one can show that
\[
R(X,Y,X,Y) = \left(\det \begin{pmatrix} X^1 & Y^1 \\ X^2 & Y^2 \end{pmatrix}\right)^2 R(e_1, e_2, e_1, e_2),
\]
and further one can easily compute that 
\[
\|X\|^2\,\|Y\|^2 - \langle X, Y \rangle^2 = \left(\det \begin{pmatrix} X^1 & Y^1 \\ X^2 & Y^2 \end{pmatrix}\right)^2. \qedhere
\]
\end{proof}






\begin{colremark}[$R$ as a quadratic form on bivectors]{}
\label{rem:R-map-on-bivectors222}
\todo[inline]{ this is kind of a duplicate of the remark \ref{rem:R-map-on-bivectors111} }

The wedge product $X \wedge Y$ here is the bivector in $\Lambda^2 T_p M$, the space of \emph{bivectors} (or 2-vectors) on $T_p M$. This is the same algebraic construction as the wedge product of covectors used to build differential forms, just applied to vectors instead. The induced inner product on $\Lambda^2 T_p M$ is
\[
\langle X \wedge Y,\, Z \wedge W\rangle = \langle X, Z\rangle \langle Y, W\rangle - \langle X, W\rangle \langle Y, Z\rangle,
\]
so in particular $\|X \wedge Y\|^2 = \|X\|^2 \|Y\|^2 - \langle X, Y\rangle^2$, matching our denominator.

By the antisymmetries of $R$ (Theorem~\ref{thm:symmetries-riemann-curvature-tensor}, items 1 and 2), the curvature tensor reduces to a bilinear map
\[
\widetilde R : \Lambda^2 T_p M \times \Lambda^2 T_p M \to \mathbb{R}, \qquad \widetilde R(X \wedge Y,\, Z \wedge W) := R(X, Y, Z, W),
\]
and the pairwise symmetry (item 3) makes $\widetilde R$ a \emph{symmetric} bilinear form on $\Lambda^2 T_p M$. In this picture the sectional curvature reads
\[
K(\Pi) = \frac{\widetilde R(X \wedge Y,\, X \wedge Y)}{\|X \wedge Y\|^2},
\]
which looks a bit like a Rayleigh-quotient: it depends only on the bivector $X \wedge Y$ up to scaling, equivalently only on the plane $\Pi = \mathrm{span}(X, Y)$.
\end{colremark}





We introduce now the  Ricci tensor. It controls volume distortion: it measures how volumes of small geodesic balls deviate from the Euclidean volume in the direction $v$. Equivalently, $\mathrm{Ric}$ governs the average behaviour of nearby geodesics: if $\mathrm{Ric}(\dot\gamma, \dot\gamma) > 0$, geodesics neighbouring $\gamma$ tend on average to converge.

\begin{colordef}{Ricci curvature}{ricci-curvature}
The \emph{Ricci tensor} $\mathrm{Ric} \in T^0_2(M)$ is the trace of the curvature operator in the first slot:
\[
\mathrm{Ric}(X, Y) := \mathrm{tr}\!\left(Z \mapsto R(Z, X)\, Y\right).
\]
\end{colordef}
In components,
\[
\mathrm{Ric}_{\alpha\beta} = R^\mu{}_{\alpha\mu\beta} = g^{\mu\nu}\, R_{\nu\alpha\mu\beta}.
\]
If $\{e_i\}_{i=1}^n$ is a local orthonormal frame, then
\[
\mathrm{Ric}(X, Y) = \sum_{i=1}^n R(e_i, X, e_i, Y).
\]

The Ricci tensor is symmetric, $\mathrm{Ric}(X, Y) = \mathrm{Ric}(Y, X)$ (see Exercis 9.1a).
\begin{colremark}[There is one unique Ricci tensor]{}
There are three possible contractions of the upper index of $R^a{}_{bcd}$ with one of the lower indices:
\begin{itemize}
    \item $R^a{}_{acd} = 0$, since this contracts a symmetric tensor ($g^{a\mu}$, used to raise the index) with one antisymmetric in $(\mu, a)$ (the antisymmetry of $R_{\mu acd}$ in its first pair).
    \item $R^a{}_{bad} = \mathrm{Ric}_{bd}$ (this is of course the Ricci tensor).
    \item $R^a{}_{bca} = -R^a{}_{bac} = -\mathrm{Ric}_{bc}$, (this is the negative Ricci tensor).
\end{itemize}
So up to sign, there is only one nontrivial contraction of $R$.
\end{colremark}






\begin{colordef}{Scalar curvature}{scalar-curvature}
The \emph{scalar curvature} $S \in C^\infty(M)$ is the trace of the Ricci tensor with respect to $g$,
\[
S := g^{ij}\, \mathrm{Ric}_{ij}.
\]
\end{colordef}
Equivalently, in a local orthonormal frame $\{e_i\}_{i=1}^n$, we have
\[
S = \sum_{i=1}^n \mathrm{Ric}(e_i, e_i) = \sum_{i,j=1}^n R(e_i, e_j, e_i, e_j).
\]

The three curvatures form a hierarchy of decreasing information:
\[
\underbrace{R}_{\text{full tensor}} \;\longrightarrow\; \underbrace{\mathrm{Ric}}_{\text{symmetric (0,2)-tensor}} \;\longrightarrow\; \underbrace{S}_{\text{scalar}}.
\]
Each is a contraction of the previous: $\mathrm{Ric}$ is a trace of $R$, and $S$ is a trace of $\mathrm{Ric}$. At each step information is lost, but the resulting object is simpler to work with.

\begin{colorlem}{Curvatures on a space of constant sectional curvature}{constant-sectional-curvature}
If $(M, g)$ has constant sectional curvature $K$, then
\[
\mathrm{Ric}_{ij} = (n-1)\, K\, g_{ij}, \quad S = n(n-1)\, K.
\]
\end{colorlem}

\begin{proof}
See handwritten notes.
\end{proof}

\begin{colexample}{Spaces of constant sectional curvature}{}
The round sphere $S^n$ (and its quotients, e.g.\ $\mathbb{RP}^n$) has constant positive sectional curvature. Hyperbolic space $\mathbb{H}^n$ has constant negative sectional curvature. Euclidean space $\mathbb{R}^n$ has $K = 0$, hence $\mathrm{Ric} = 0$ and $S = 0$.
\end{colexample}












\subsection{Riemannian submanifolds}

We introduce some natural notions for submanifolds, and characterize the geometry of submanifolds in terms of the curvature tensors.

Let $N^n$ be a submanifold of $(M^m,g)$, so locally, for all $p\in N$, there exists a coordinate system $(x^1, \ldots, x^m)$ on $U\subset M$ such that 
$$N\cap U=\{p\in M: x^{n+1}(p)=....=x^m(p)=0\}.$$

The tangent space of $N$ at $p$ is 
$$T_pN = \{v\in T_pM: v^{n+1} = \cdots = v^m = 0\},$$

where $v^i$ are the coordinates of $v$ in the basis $\partial_i$. There is also the \emph{perpendicular space} to $N$ at $p$:
$$T_pN^\perp = \{w\in T_pM: w^1 = \cdots = w^n = 0\}=\{w\in T_pM: w\perp T_p N\}$$

where $\perp$ is to be understood with respect to the metric $g$. We define the orthogonal projections 
\begin{align}
\label{eq:projections_submanifold1}
\pi^T: T_pM &\to T_pN \\ 
\label{eq:projections_submanifold2}
\pi^\perp: T_pM &\to T_pN^\perp.
\end{align}

So all $X$ can be decomposed as 
$$v = v^T + v^\perp,$$

where $v^T = \pi^T(v)$ and $v^\perp = \pi^\perp(v)$, because $T_pM = T_pN \oplus T_pN^\perp$.

\begin{colordef}{Metric on Riemannian submanifold}{}
We define the \emph{induced metric} (first fundamental form) on $N$ by $\bar g := g|_{T_pN}$, i.e. at each $p\in N$, we take the bilinear form $g_p: T_pM \times T_pM \to \mathbb{R}$ and restrict both arguments to vectors in $T_pN$:

\[
\bar{g}_p(u,v) = g_p(u,v), \quad u,v \in T_pN.
\]
\end{colordef}



For any vector fields $X, Y$ on $N$, we decompose the covariant derivative (computed in $M$) into tangential and normal components:
$$\nabla_X Y = \pi^T(\nabla_X Y) + \pi^\perp(\nabla_X Y),$$
where the first term is the induced connection on $N$, and the second one the \emph{second fundamental form}, also denoted by $B(X,Y)=\pi^\perp(\nabla_X Y)\in TN^\perp$.


